\documentclass[12pt]{article}
\usepackage[utf8]{inputenc}
\usepackage[T1]{fontenc}
\usepackage{hyperref}
\title{Yoruba Language Resources for K-12 Curriculum Development}
\author{Compiled by Claude Sonnet 5.5 (claude-sonnet-5-5) for Andr\'as Kornai}
\date{}
\begin{document}
\maketitle

\section{Introduction}
Yoruba (ISO 639-3 \texttt{yor}, Glottocode yoru1245) has about 40 million speakers, mainly in south-western Nigeria. Nigeria's National Policy on Education calls for the first three years of primary school to be taught in the mother tongue. This report is based on a single web search, so it is a starting point and not an exhaustive survey. Coverage of dictionaries, grammars and Bible translations is missing.

\section{Primary Educational Resources}
\begin{itemize}
\item USAID donated over 10,000 Yoruba-language textbooks for Primary 1--3 to pupils in Oyo State, developed with language experts and the Oyo State UBE Board: \url{https://nannews.ng/usaid-donates-textbook-in-local-language-to-pupils-in-oyo/}
\item A commercial site sells a Primary 1 Yoruba scheme of work for the 2026 revised national curriculum. I could not verify it against an official source: \url{https://www.headteacher.ng/product/primary-1-yoruba-language-scheme-of-work-2026-new-national-curriculum/}
\end{itemize}

\section{Linguistic and Reference Materials}
The search found research literature on Yoruba in education:
\begin{itemize}
\item Quality control in Yoruba language textbooks for primary schools in Nigeria: \url{https://journals.ezenwaohaetorc.org/index.php/IJAAS/article/view/1-2016-010}
\item Teachers' experiences of using Yoruba as a medium of instruction in primary classes: \url{https://repository.up.ac.za/items/a8bd99a3-8158-47eb-a270-7589cb10ee3c/full}
\item University of Ibadan repository item on Yoruba in lower primary classes: \url{https://repository.ui.edu.ng/items/9a7877d6-9817-4355-9aff-332ba6ec28db}
\end{itemize}

\section{Religious and Community Resources}
None were found in this search.

\section{Digital and Audio Resources}
None were found in this search.

\section{Curriculum Development Context}
Studies of lower primary classes in Ibadan report shortages of readers and textbooks, a need for trained Yoruba teachers, and weak monitoring by local authorities. These gaps are the main opportunity for new materials.
\end{document}